\documentclass[11pt]{article}
\usepackage[T1]{fontenc}
\usepackage[utf8]{inputenc}
\usepackage{lmodern}
\usepackage{amsmath,amssymb,amsthm,mathtools}
\usepackage[margin=1in]{geometry}
\usepackage{longtable,booktabs,array}
\usepackage{xurl}
\usepackage[hidelinks]{hyperref}
\newtheorem{theorem}{Theorem}[section]
\newtheorem{proposition}[theorem]{Proposition}
\newtheorem{lemma}[theorem]{Lemma}
\newtheorem{corollary}[theorem]{Corollary}
\theoremstyle{remark}
\newtheorem{remark}[theorem]{Remark}
\DeclareMathOperator{\Gal}{Gal}
\allowdisplaybreaks
\setlength{\emergencystretch}{3em}
\title{ZETA\_NEXT\\Effective class numbers from quadratic zero exclusion,\\
bad fibres of $\Lambda$-schemes, and the non-CM gate}
\author{GPT-6 Astra Pro\\Research run of 21 November 2026}
\date{21 November 2026}
\begin{document}
\maketitle

\begin{abstract}
Using the supplied theorem of Alp\"oge and the explicitly checked
constants in Stark's exceptional-zero argument, we prove
$\operatorname{Res}_{s=1}\zeta_K(s)>1/(173\log D_K)$ for every normal
$K\ne\mathbb Q$, and $1/(173n!\log D_K)$ for every field of degree
$n\ge2$.  This gives effective Brauer--Siegel inequalities, explicit
fixed-degree CM class-number bounds, and explicit Galois-orbit bounds
for every CM point of $\mathcal A_g$, including products and
nonmaximal endomorphism orders.  The orbit bound has exponent
$1/(28672g^4)$ and a completely specified positive constant.
Only the averaged Colmez formula is needed.

The proposed bad-fibre statement for flat finite-type $\Lambda$-schemes
is false without separatedness: a gluing construction inserts the
zeta function of an arbitrary affine finite-field scheme into a single
bad factor.  For proper flat $\Lambda$-schemes the statement is true,
on actual cohomology: a Frobenius-tower comparison identifies special
cohomology with a Galois-stable summand of generic cohomology.
We also prove the corrected proper $\Lambda$-RH conclusion.

For absolute-Hodge motives of abelian type we construct the canonical
maximal CM summand as the invariants of the connected derived motivic
group.  We give exact residual $L$-function lists for curves and
abelian varieties, explicitly evaluate the full symplectic case and
powers of a non-CM elliptic curve, and exhibit a simple non-CM
Jacobian with a rank-two non-Tate CM summand in primitive degree four.
The splitting theorems are unconditional.  Applying RH to their CM
factors remains conditional on the stated Masullo claim; RH for their
non-CM factors is identified as the remaining analytic input.
\end{abstract}

\tableofcontents
\clearpage
\section*{Reading conventions, source status, and the answer to the task}
\addcontentsline{toc}{section}{Reading conventions and answer to the task}

This paper continues ZETA\_TWO and its intake of 21 November 2026.
Its IA01 audit and \path{CONDITIONAL_AUDIT.tsv} stand as requested;
we do not repeat them.  A \emph{finding} means a proved deduction or
counterexample established in this work.  It does not assert novelty
against all of the literature.  Every finding below has received a
complete proof reading; the methods, actors, dates, and limited scope
of each check appear in \path{CHECK_EVENTS.tsv}.

\paragraph{The supplied zero theorem A.}
We cite as a theorem the supplied four-page version of L.~Alp\"oge,
\emph{Remark} \cite{Alpoge}.  The endpoint used here is step~20 on
p.~4, not step~17.  In \path{ALPOGE_Main.lean} the declarations are
\path{riemannHypothesis}, \path{riemannZeta_eq_zero_re},
\path{LFunction_ne_zero_of_half_lt_re},
\path{LFunction_eq_zero_re}, and
\path{LFunction_eq_zero_re_of_isPrimitive}; their precise predicates
are in \path{ALPOGE_Statements.lean}.  We use RH for $\zeta$ and
Dirichlet GRH, especially nonvanishing to the right of $1/2$ for
quadratic characters.  We do not infer GRH for a general Dedekind
zeta function from this input.  No step of the theta, positivity,
or Lean argument is re-audited.

\paragraph{The separate conditional input M.}
The Masullo input is the draft's claim in \S1.2, PDF p.~12, for all
unitary Hecke characters of every number field, including
infinite-order characters.  That assertion is unnumbered in the
supplied draft.  Theorem~13.3.1.2, PDF p.~522, is its intended
\emph{Divisor of $L(s,\chi)$} endpoint but has no completed statement
or proof; \S14.1.2, PDF p.~528, discusses the intended consequence.
These are the exact source statuses already recorded in ZETA\_TWO
MDOC01 \cite{Masullo,ZetaTwo}.  Every use of M in this paper is
explicitly conditional.  Anderson's same-weight factorization,
Corollary~5.7.2 and Remark~5.7.3, p.~181 \cite{Anderson1986}, is
carried through ZETA\_TWO CM01; no new proof of that known theorem
is claimed.

\paragraph{Principal answers.}
For a number field write $D_K=|\operatorname{disc}K|$,
$w_K=|\mu(K)|$ and $(r_1,r_2)$ for its signature.  NX01--NX03 prove
\[
 h_KR_K>
 \frac{w_K\sqrt{D_K}}
 {173\,2^{r_1}(2\pi)^{r_2}\log D_K}
 \quad(K/\mathbb Q\ \hbox{normal},\ K\ne\mathbb Q),
\]
and give the explicit factorial loss for arbitrary degree.
NC01--NC04 propagate this to CM class numbers, heights, and all CM
orbits.  In particular, for an explicit $c_g>0$ and all CM fields
$E$ of degree $2g$,
\[
 h_E>c_gD_E^{1/4}(\log D_E)^{-g}.
\]
The full orbit formula, with its constants defined before and in
NC04, is
\[
 |\Gal(\overline{\mathbb Q}/\mathbb Q)\cdot x|
 \ge \frac{|\operatorname{disc}Z(\operatorname{End}_{\overline{\mathbb Q}}A_x)|^{1/P_g}}
 {s_gT_g^{1/P_g}},\qquad P_g=28672g^4.
\]
There is no residual ineffective constant in these displayed
statements.  They do not themselves supply an algorithm for every
other step of a general Andr\'e--Oort proof.

NX-L01 refutes the exact bad-fibre conjecture asked in the prompt:
its hypotheses do not include separatedness.  NX-L02--NX-L03 prove
the proper-flat splitting assertion, and NX-L04 proves that the
multiplicities need not be the generic inertia multiplicities.
NX-L05 explains a failed proper counterexample route; NX-L06 gives
meromorphic continuation and half-integer divisor lines in the
proper case using A.

NG01 closes ZT06's CM-summand gap for every pure factor of abelian
type.  NG02--NG03 compute the residual lists; NG04--NG07 make their
content concrete.  For a general abelian variety the list is an
exact abstract formula; evaluating its primitive CM projectors
requires the actual Mumford--Tate representation.  No general
algorithm recovering that representation from equations is asserted.
The non-CM RH statements which the question asks us to name remain
the analytic inputs explicitly identified by these formulas.

\paragraph{Finite factors and the meaning of $\zeta_X$.}
For a variety over $\mathbb Q$, the Euler product is understood
through a chosen model $\mathcal X$ over an open of
$\operatorname{Spec}\mathbb Z$; its incomplete form is denoted by
$\zeta_X^S$.  For a finite-type scheme over $\mathbb Z$ it is the
closed-point Euler product.  Restoring bad primes always uses the
actual fibres, as in ZETA\_TWO ZT04.  No bad local factor is silently
replaced by the canonical generic inertia factor.

\paragraph{Archive conventions.}
The available cote-and-scan locators are carried from the supplied
predecessor and intakes, not newly viewed manuscripts.  They are
listed precisely in the final framework table and
\path{IMPLICATIONS.md}.  None of the new analytic estimates,
perfectoid comparisons, or Faltings--Deligne theorems is attributed
to an unavailable Grothendieck scan.  Where the supplied material
does not provide a cote and scan, that absence is stated.

\part{Checked findings}
